\chapter{اللقاء الحادي والثلاثون: كتمان ما في الأرحام والطلاق مرتان والخلع}

\section{مقدمة}
السلام عليكم ورحمة الله وبركاته. الحمد لله رب العالمين، وأشهد أن لا إله إلا الله وحده لا شريك له، وأشهد أن محمداً عبده ورسوله، صلى الله عليه وعلى آله وصحبه وسلم.

استأنف الشيخ هذا اللقاء من الموضع الذي توقف عنده في المجلس السابق، فدخل في بقيّة أحكام العدة، ثم انتقل إلى أوائل أحكام الطلاق والخلع، من قوله تعالى: \begin{ayah}وَلَا يَحِلُّ لَهُنَّ أَنْ يَكْتُمْنَ مَا خَلَقَ اللَّهُ فِي أَرْحَامِهِنَّ\end{ayah} إلى قوله تعالى: \begin{ayah}تِلْكَ حُدُودُ اللَّهِ فَلَا تَعْتَدُوهَا\end{ayah}.

وفي هذا اللقاء ظهر صنيع \rawi{الطبري} في غاية الوضوح: جمع الأقوال، وبيان سبب كل خلاف، والنظر في السياق، ثم رد الفروع إلى الأصول الكبرى في باب الأسرة والحقوق والحدود.

\separator

\section{تأويل (ولا يحل لهن أن يكتمن ما خلق الله في أرحامهن)}
\isnad{قال أبو جعفر:} اختلف أهل التأويل في المراد بما خلق الله في الأرحام، فحكى \rawi{الطبري} أقوالهم على أنحاء:
\begin{itemize}
    \item فقيل: المراد الحيض.
    \item وقيل: المراد الحمل.
    \item وقيل: المراد الأمران جميعاً: الحيض والحمل.
\end{itemize}

ثم رجح أن الآية تتناول الأمرين معاً، لأن كلاًّ منهما يترتب عليه حكم الرجعة وانقضاء العدة. فالمرأة إذا كتمت الحيض أو الحمل أبطلت بذلك حقاً أثبته الله لزوجها في زمن العدة، أو أوقعت اضطراباً في الأحكام المترتبة على الطلاق.

وبيّن الشيخ أن من محاسن ترجيح \rawi{الطبري} هنا أنه لم يقف عند مجرد الاحتمال اللغوي، بل نظر إلى أثر الكتمان في الحكم الشرعي، فكل ما يترتب على إخفائه إبطال حق الزوج أو خلط أمر العدة دخل في النهي.

\begin{fawaid}[الترجيح بالأثر العملي للحكم من وجوه القوة]
من دقة \rawi{الطبري} أنه لا يوازن بين الأقوال من جهة اللفظ وحده، بل ينظر أيضاً إلى ما يترتب على كل قول من الأحكام. فحيث اشترك الحيض والحمل في التأثير على العدة والرجعة، كان إدخالهما في الآية أظهر وأقوى.
\end{fawaid}

\section{تأويل (وبعولتهن أحق بردهن في ذلك)}
قرر \rawi{الطبري} أن هذا الحكم متعلق بالمطلقة الرجعية، أي من طلقت طلقة أو طلقتين، ما دامت في عدتها. فالبعل أحق بردها في تلك المدة إذا أراد الإصلاح، لأن الله جعل له هذا الحق ما لم تنقض العدة أو تقع البينونة التامة.

ومن هنا شدد الشيخ على أن إخفاء المرأة ما في رحمها ليس مجرد مخالفة خفية، بل قد يكون سبباً مباشراً في إسقاط حق ثابت، أو في إلحاق الولد بغير موضعه، أو في التحايل على الفكاك من حكم الرجعة قبل تمام أجله.

\separator

\section{تأويل (ولهن مثل الذي عليهن بالمعروف وللرجال عليهن درجة)}
وقف الشيخ مع هذا الموضع وقفة تربوية نافعة، فبيّن أن الشريعة لا تقيم العلاقة الزوجية على منطق المغالبة، بل على العدل، وأداء كل طرف ما عليه، ومراعاة ما جعله الله له من الحق.

ثم نبّه إلى أن \rawi{الطبري} يربط هذه الأحكام بسياقها، فلا تُقرأ آيات الأسرة قراءة مبتورة، بل تُفهم ضمن منظومة من الحقوق، والرجعة، والعشرة، والعدل، وترك الظلم.

\section{تفسير خاتمة الآية: (والله عزيز حكيم)}
فسر \rawi{الطبري} ختم هذا الموضع باسمَي الله \textit{العزيز الحكيم} تفسيراً مناسباً لسياق الأحكام السابقة:
\begin{itemize}
    \item فهو \textit{عزيز} في انتقامه ممن خالف أمره وتعدى حدوده في هذه الأبواب.
    \item وهو \textit{حكيم} فيما دبره وقضاه من تشريعات الأسرة والطلاق والرجعة.
\end{itemize}

وهذا من بديع التفسير؛ لأن العزة هنا تردع من يستهين بهذه الأحكام، والحكمة تطمئن القلب إلى أن هذه الحدود ليست عشوائية، بل موضوعة في غاية العدل والإحكام.

\begin{fawaid}[ربط الأسماء الإلهية بالسياق يزيد الحكم وضوحاً]
إذا ختمت الآية بـ \textit{العزيز الحكيم} بعد جملة من أحكام الأسرة، دل ذلك على أن التشريع مؤيد بالقهر لمن خالفه، وبالحكمة لمن تدبره، لا أنه مجرد أوامر منفصلة عن صفات الرب سبحانه.
\end{fawaid}

\separator

\section{تأويل (الطلاق مرتان)}
ذكر \rawi{الطبري} في هذه الآية قولين مشهورين:
\begin{itemize}
    \item قيل: هي بيان لعدد الطلاق الذي يملك فيه الزوج الرجعة، وأنه طلقتان، ثم إذا أوقع الثالثة بانت المرأة منه.
    \item وقيل: هي بيان لسنة الطلاق وكيفية إيقاعه في الأطهار.
\end{itemize}

ثم رجح أن المقصود الأول هو الأظهر: أن الآية مسوقة لبيان عدد الطلاق الذي تكون معه الرجعة، وأن التحريم التام إنما يكون بعد الثالثة. واستند في ذلك إلى سياق الآيات التي بعدها، وإلى ما ورد من الخبر في تفسير قوله: \begin{ayah}فَإِمْسَاكٌ بِمَعْرُوفٍ أَوْ تَسْرِيحٌ بِإِحْسَانٍ\end{ayah}.

وهذا الترجيح مثال واضح على أن السياق من أعظم المرجحات؛ فالآية ليست مجرد حكم مفرد، بل جزء من نسق متتابع يبين متى تبقى الرجعة، ومتى تنقطع، ومتى تحرم المرأة حتى تنكح زوجاً غيره.

\begin{fawaid}[السياق من أقوى المرجحات بين الأقوال]
قد يكون لكل قول حظ من النظر، لكن الذي يرده \rawi{الطبري} كثيراً إلى الصواب هو النظر في انتظام الآية مع ما قبلها وما بعدها. فالسياق هنا دل على أن المقصود الأول بيان عدد الطلاق الذي يملك فيه الزوج الرجعة.
\end{fawaid}

\section{تأويل (فإمساك بمعروف أو تسريح بإحسان)}
في هذه الجملة نقل \rawi{الطبري} خلافاً آخر:
\begin{itemize}
    \item فقيل: المراد بالتسريح هنا التطليقة الثالثة.
    \item وقيل: بل المراد ترك المرأة حتى تنقضي عدتها من غير مراجعة.
\end{itemize}

ثم مال إلى أن التسريح بالإحسان هو التطليقة الثالثة، بناء على الخبر المروي في ذلك، وأن الإمساك بالمعروف هو مراجعتها مع إحسان العشرة وكف الظلم، لا مجرد إرجاعها صورياً مع بقاء الإضرار.

وشرح الشيخ أن هذا من المواضع التي يتقدم فيها الخبر إذا صح عند \rawi{الطبري} على مجرد الاحتمال اللغوي، لأن النص المرفوع يرفع التردد ويعين المراد.

\begin{fawaid}[إذا صح الخبر ارتفع كثير من الاحتمال]
من منهج \rawi{الطبري} أن الاحتمال الذي يحتمله ظاهر اللفظ لا يبقى على إطلاقه إذا جاء الخبر المبين عن رسول الله صلى الله عليه وسلم، بل يصير هو المقدم في تعيين المراد.
\end{fawaid}

\separator

\section{تأويل (ولا يحل لكم أن تأخذوا مما آتيتموهن شيئاً)}
قرر \rawi{الطبري} أن الأصل في هذا الباب منع الأزواج من أخذ شيء مما آتوا نساءهم عند الطلاق، لأن العشرة لا يجوز أن تتحول إلى باب أكل للحقوق أو تعسف في الفراق.

ثم جاء الاستثناء بعد ذلك في باب الخلع، فكان من المهم أن يميز القارئ بين الأصل العام، وبين الصورة الخاصة التي أبيح فيها الأخذ على وجه مخصوص.

\section{تأويل (إلا أن يخافا ألا يقيما حدود الله)}
هنا دخل \rawi{الطبري} في تحرير باب الخلع، فذكر أقوالاً كثيرة في معنى الخوف المذكور:
\begin{itemize}
    \item فقيل: يكون إذا ظهر النشوز من المرأة خاصة.
    \item وقيل: يكون إذا صرحت ببغضها وعدم قدرتها على القيام بحقوق الزوج.
    \item وقيل: بل المعتبر خوف الوقوع في معصية الله من الجانبين جميعاً بسبب فساد العشرة بينهما.
\end{itemize}

ثم رجح أن المعتبر هو خوف المسلمين عليهما جميعاً ألا يقيما حدود الله، أي: أن الحال بين الزوجين قد آل إلى صورة يخشى معها أن يفرط كل واحد منهما في حق صاحبه. وهذا لا ينافي أن يكون مبدأ الفساد من جهة المرأة؛ لأن نشوزها قد يجر الزوج إلى التقصير أو سوء المجازاة، فيتحقق الخوف من الجانبين.

وبيّن الشيخ أن هذا من أدق مواضع الكتاب في تحرير محل النزاع؛ إذ ليس المراد مجرد وجود بغض نفسي عابر، ولا أن يضار الرجل زوجته حتى تفتدي منه، بل المراد صورة تتعذر معها إقامة حدود الله في الصحبة والعشرة.

\begin{fawaid}[اختلاف الصور يوجب اختلاف الأحكام]
ليس كل أخذ من مال الزوجة في الفراق سواء. فهناك صورة محرمة إذا كان الزوج هو المضار، وصورة مستثناة إذا خيف من الجانبين ألا تقام حدود الله. وبهذا يظهر أن اختلاف الأحكام تابع لاختلاف الصور، لا أن النصوص متعارضة في نفسها.
\end{fawaid}

\section{فائدة منهجية في باب النسخ والجمع}
أفاد الشيخ هنا فائدة جليلة في منهج النظر: أن اختلاف الحكم لا يقتضي النسخ إلا إذا اتحدت الجهة والمحل وتعذر الجمع. أما إذا اختلفت الصورة، أو اختلف محل الحكم، أمكن الجمع، وكان ذلك من اختلاف التنوع أو من اختلاف الأحكام باختلاف الأحوال.

ولهذا لم يجعل \rawi{الطبري} آية منع الأخذ من النساء منسوخة بآية الخلع، بل جعل كل واحدة منهما نازلة على صورة غير صورة الأخرى، فزال التعارض، وبقيت الآيتان محكمتين في موضعهما.

\begin{fawaid}[لا يثبت النسخ مع إمكان حمل كل نص على حاله]
من القواعد النفيسة أن النصين إذا أمكن تنزيل كل واحد منهما على صورة غير صورة الآخر، لم يجز التسرع إلى دعوى النسخ. وهذه قاعدة مهمة جداً في أبواب التفسير والفقه جميعاً.
\end{fawaid}

\separator

\section{تأويل (تلك حدود الله فلا تعتدوها)}
ختم \rawi{الطبري} هذا المقطع بأن هذه الأحكام كلها من حدود الله، أي: من معالم ما أحل وما حرم، وما أمر به ونهى عنه. فليس الكلام عن الأسرة والطلاق والرجعة والخلع شأناً عرفياً محضاً، بل هو باب تعبدي تحكمه حدود الله.

ولهذا نبّه الشيخ إلى فائدة عظيمة: أن تقوى الله لا تكون بالتمني، بل بمعرفة حدوده أولاً، ثم الوقوف عندها ثانياً. فالجهل بالحدود سبب من أسباب التعدي، كما أن العمد إلى تجاوزها سبب آخر للظلم.

\begin{fawaid}[لا تتحقق التقوى إلا بعد معرفة الحدود]
من أراد أن يتقي الله في باب من الأبواب فلابد أن يتعلم أولاً حدود الله فيه. فالزوج والزوجة، والتاجر، والحاج، وكل مكلف، لا يعذر نفسه بجهله ثم يظن أنه متق لله.
\end{fawaid}

\separator

خرج هذا اللقاء بأصلين كبيرين: أن أحكام الأسرة لا تفهم إلا بردها إلى حدود الله، وأن كثيراً من الإشكالات في الطلاق والخلع والرجعة تزول إذا أحسن القارئ تصور الصور المختلفة التي نزلت عليها الأحكام، ولم يخلط بين مواضع المنع ومواضع الاستثناء.
